\chapter{Expansion in the hypercube graph}
\label{chap:expansion}

\section{The graph and its spectrum}


\begin{definition}[Rerandomizing one coordinate]\label{def:rerandomize-coord}
  \lean{MIPStarRE.Paper2009LDT.Section7ExpansionHypercubeGraph.rerandomizeCoord}
  For $u \in \F_q^m$, an index $i \in \{1,\dots,m\}$, and $x \in \F_q$, let
  \[
    \operatorname{rerand}_i(u,x) = u + x e_i.
  \]
  This is the move obtained by rerandomizing the $i$-th coordinate of $u$ by a uniform additive shift.
\end{definition}

\begin{definition}[Hypercube graph]\label{def:hypercube-graph}
  \uses{def:rerandomize-coord}
  The hypercube graph $C=(V,E)$ has vertex set $V=\F_q^m$, and an edge between $u$ and $v$ whenever they differ in at most one coordinate. A random edge $(u,v) \sim C$ is sampled by drawing $u \sim \F_q^m$, $i \in \{1,\dots,m\}$, and $x \in \F_q$ uniformly and then setting $v=\operatorname{rerand}_i(u,x)$.
\end{definition}

\begin{definition}[Adjacency matrix and Laplacian]\label{def:adjacency-laplacian}
  \uses{def:hypercube-graph}
  Let $M=q^m$. The normalized adjacency matrix of $C$ is
  \[
    K = \E_{(u,v)\sim C} \ket{u}\bra{v},
  \]
  and the Laplacian is
  \[
    L = \frac{1}{M}I-K.
  \]
\end{definition}

\begin{lemma}[Laplacian rewrite]\label{prop:laplacian-rewrite}
  \lean{MIPStarRE.Paper2009LDT.Section7ExpansionHypercubeGraph.laplacianRewrite}
  \uses{def:adjacency-laplacian}
  The Laplacian can be written as
  \[
    L = \frac12 \E_{(u,v)\sim C} (\ket{u}-\ket{v})(\bra{u}-\bra{v}).
  \]
\end{lemma}
\begin{proof}
  The two vertex marginals of a random edge are uniform on $\F_q^m$. Expanding the right-hand side yields $(1/M)I-K$.
\end{proof}

\begin{lemma}[Average of a nontrivial additive character]\label{lem:character-average-scalar}
  For $\beta \in \F_q$,
  \[
    \E_{x \sim \F_q} \omega^{\operatorname{tr}(\beta x)} =
    \begin{cases}
      1 & \text{if } \beta=0,\\
      0 & \text{if } \beta\ne 0.
    \end{cases}
  \]
\end{lemma}
\begin{proof}
  If $\beta=0$ the character is constant. If $\beta\ne 0$, the map $x \mapsto \beta x$ permutes $\F_q$, so the average is the average of a nontrivial additive character, which vanishes.
\end{proof}

\begin{lemma}[Orthogonality of finite-field characters]\label{lem:character-average-vector}
  \uses{lem:character-average-scalar}
  For $\alpha,\beta \in \F_q^m$,
  \[
    \frac{1}{q^m}\sum_{u \in \F_q^m} \omega^{\operatorname{tr}(u\cdot(\beta-\alpha))} =
    \begin{cases}
      1 & \text{if } \alpha=\beta,\\
      0 & \text{if } \alpha\ne \beta.
    \end{cases}
  \]
\end{lemma}
\begin{proof}
  Factor the sum over the $m$ coordinates and apply Lemma~\ref{lem:character-average-scalar} to each factor.
\end{proof}

\begin{lemma}[Eigenvectors of the hypercube graph]\label{prop:eigenvectors}
  \lean{MIPStarRE.Paper2009LDT.Section7ExpansionHypercubeGraph.eigenvectors}
  \uses{def:adjacency-laplacian, lem:character-average-vector, lem:character-average-scalar}
  For each $\alpha \in \F_q^m$, define
  \[
    \ket{\varphi_\alpha} = \frac{1}{q^{m/2}} \sum_{u \in \F_q^m} \omega^{\operatorname{tr}(u\cdot \alpha)} \ket{u}.
  \]
  Then the vectors $\ket{\varphi_\alpha}$ form an orthonormal basis of $\C^V$, and $\ket{\varphi_\alpha}$ is an eigenvector of $K$ with eigenvalue
  \[
    \frac{1}{q^m}\cdot \frac{m-|\alpha|}{m},
  \]
  where $|\alpha|$ is the number of nonzero coordinates of $\alpha$.
\end{lemma}
\begin{proof}
  The orthonormality is exactly Lemma~\ref{lem:character-average-vector}. For the eigenvalue, substitute the edge-sampling rule $v=u+x e_i$ into $K\ket{\varphi_\alpha}$, separate the average over $u$ from the average over the rerandomized coordinate, and use Lemma~\ref{lem:character-average-scalar} to replace $\E_x\,\omega^{\operatorname{tr}(x\alpha_i)}$ by $\mathbf 1_{\alpha_i=0}$. Averaging over $i$ leaves the factor $(m-|\alpha|)/m$.
\end{proof}

\begin{lemma}[Spectral gap of the Laplacian]\label{cor:laplacian-spectral-gap}
  \lean{MIPStarRE.Paper2009LDT.Section7ExpansionHypercubeGraph.laplacianSpectralGap}
  \uses{prop:eigenvectors}
  If $\lambda_1 \le \lambda_2 \le \cdots \le \lambda_{q^m}$ are the eigenvalues of $L$, then
  \[
    \lambda_1=0,
    \qquad
    \lambda_2 = \frac{1}{m q^m}.
  \]
\end{lemma}
\begin{proof}
  Lemma~\ref{prop:eigenvectors} identifies the two largest eigenvalues of $K$, hence the two smallest eigenvalues of $L=(1/q^m)I-K$.
\end{proof}

\section{Local and global variance}

\begin{definition}[Local and global variance]\label{def:local-and-variance}
  \lean{MIPStarRE.Paper2009LDT.Section7ExpansionHypercubeGraph.localAndVariance}
  Let $\ket{\psi} \in \mathcal H_{\mathrm A} \ot \mathcal H_{\mathrm B}$ and let $0 \le A^u \le I$ be an operator for each $u \in \F_q^m$. The local variance is
  \[
    \mathbf{Var}_{\mathrm{local}}(A,\psi)
    = \frac12 \E_{(u,v)\sim C} \bra{\psi} (A^u-A^v)^2 \ot I \ket{\psi},
  \]
  and the global variance is
  \[
    \mathbf{Var}_{\mathrm{global}}(A,\psi)
    = \frac12 \E_{u,v\sim \F_q^m} \bra{\psi} (A^u-A^v)^2 \ot I \ket{\psi}.
  \]
\end{definition}

\begin{lemma}[Local variance as a Laplacian form]\label{lem:local-rewrite}
  \lean{MIPStarRE.Paper2009LDT.Section7ExpansionHypercubeGraph.localRewrite}
  \uses{def:adjacency-laplacian, def:local-and-variance}
  Define
  \[
    A_{\mathrm{comb}} = \sum_{u \in \F_q^m} \ket{u} \ot A^u \ot I.
  \]
  Then
  \[
    \Tr(A_{\mathrm{comb}}^\dagger (L \ot \ket{\psi}\bra{\psi}) A_{\mathrm{comb}})
    = \mathbf{Var}_{\mathrm{local}}(A,\psi).
  \]
\end{lemma}
\begin{proof}
  Insert the expression for $L$ from Lemma~\ref{prop:laplacian-rewrite} and evaluate the resulting quadratic form on the vectors $\ket{u}-\ket{v}$.
\end{proof}

\begin{lemma}[Global variance as the orthogonal Fourier mass]\label{lem:global-rewrite}
  \lean{MIPStarRE.Paper2009LDT.Section7ExpansionHypercubeGraph.globalRewrite}
  \uses{prop:eigenvectors, def:local-and-variance}
  Write
  \[
    A_{\mathrm{comb}} = \ket{\varphi_0} \ot A_0 + \ket{\varphi_\perp} \ot A_\perp,
  \]
  where $\ket{\varphi_\perp}$ is orthogonal to $\ket{\varphi_0}$. Then
  \[
    \frac{1}{q^m}\Tr((\bra{\varphi_\perp}\ot A_\perp)(I\ot \ket{\psi}\bra{\psi})(\ket{\varphi_\perp}\ot A_\perp))
    = \mathbf{Var}_{\mathrm{global}}(A,\psi).
  \]
\end{lemma}
\begin{proof}
  The $\ket{\varphi_0}$ component is the average operator $A_{\mathrm{avg}}$. Subtracting it from $A_{\mathrm{comb}}$ leaves the fluctuations $A^u-A_{\mathrm{avg}}$, whose square average is the global variance.
\end{proof}

\begin{lemma}[Local-to-global inequality]\label{lem:local-to-global}
  \lean{MIPStarRE.Paper2009LDT.Section7ExpansionHypercubeGraph.localToGlobal}
  \uses{cor:laplacian-spectral-gap, lem:local-rewrite, lem:global-rewrite}
  For every family $\{A^u\}$,
  \[
    \mathbf{Var}_{\mathrm{global}}(A,\psi) \le m\, \mathbf{Var}_{\mathrm{local}}(A,\psi).
  \]
\end{lemma}
\begin{proof}
  The zero Fourier mode does not contribute to the local form. On the orthogonal complement, the Laplacian is bounded below by its spectral gap from Lemma~\ref{cor:laplacian-spectral-gap}. Compare the two quadratic forms using Lemmas~\ref{lem:local-rewrite} and~\ref{lem:global-rewrite}.
\end{proof}
